\documentclass[t,aspectratio=169]{beamer}

\usepackage{soul}
\usepackage{tikz}
\usepackage{multicol}
\usepackage[utf8]{inputenc}

% ----------------------------------------------------------------------------------------------------------------------------------------

\title{\includegraphics[width=45pt]{oepecsetlogo.png}\\(KTXFI1EBNF) Physics I. Practice}

\author{Dr.~Gábor FACSKÓ, PhD\\\tiny Assistant Professor / Senior Research Fellow\\\textit{facsko.gabor@uni-obuda.hu}}

\institute{\Tiny Óbuda University, Faculty of Electrical Engineering, 1084 Budapest, Tavaszmező u. 17.\\Wigner Research Centre for Physics, Department of Space Physics and Space Technology, 1121 Budapest, Konkoly-Thege Miklós út 29–33.\\ \url{https://wigner.hu/~facsko.gabor}}

\date{February 23,~2026}

\logo{\includegraphics[height=0.9 true cm]{kekcsik.png}}

% ----------------------------------------------------------------------------------------------------------------------------------------

\begin{document}

\frame{\titlepage}

% ----------------------------------------------------------------------------------------------------------------------------------------

\begin{frame}[fragile,squeeze,allowframebreaks]
\frametitle{Exercises}
\begin{itemize}
\setlength{\parskip}{0pt}
\setlength{\itemsep}{0pt}
\item \textbf{Mechanics of Rigid Bodies I.} Concept of a rigid body. Forces acting on a rigid body. Motion of a rigid body. Concept of the centre of mass. Location of the centre of mass. Torque. Angular momentum. Energy of a rotating rigid body, angular momentum of a rotating rigid body, Steiner’s theorem.
\item \textbf{Mechanics of Rigid Bodies II.} Principal theorems of rigid body motion; conservation laws. Collisions.
\item \textbf{Oscillations I.} Dynamic condition of harmonic oscillatory motion. Differential equation of harmonic undamped oscillatory motion.
\item \textbf{Oscillations II.} Superposition of oscillations. Decomposition of oscillations into components. Damped oscillations. Forced oscillations, resonance. Coupled oscillations.
\item Example~1: Projectile motion: \textit{Find the launch angle that maximizes the range of a projectile for a given initial speed $v$.}

\pagebreak
\item Example 2: Projectile motion: \textit{Calculate the impact distance of a projectile given an initial speed $v=1500\,km/s$ and an angle $\alpha=75^{\circ}$. At what other angle is the range identical?}
\item Non-conservative forces, Friction
\end{itemize}
\end{frame}


% ----------------------------------------------------------------------------------------------------------------------------------------

\begin{frame}
%\frametitle{Minden jó, ha a vége jó}

\vfill
\begin{center}
\textbf{\huge The End}

\bigskip
Thank you for your attention!
\end{center}
\vfill
\textbf{Acknowledgements} Thank you for the course materials for Szilárd~ZSÓKA and Zoltán~VARGA.
\end{frame}

\end{document}
